% =============================================================================
% Appendix D — Mathematical Theorems (9 R18 theorems + 3 appendix-only
% contributions T6 + T8 + T9)
%
% Status: SHIPPED (Wf-Math-1, 2026-09-26).
%
% This appendix consolidates the nine mathematical theorems (T1 -- T9) that
% underpin the R18 upgrade path of Mol-Metal.  Six of the nine (T1, T2, T5
% shipped as math + adapter; T3, T4, T7 shipped as math only) are referenced
% from the methodology sections (§3.4 SE(3) SB, §3.5 Transport Relaxation,
% §3.6 Multi-Parameter PH, §3.7 Monge-Ampere / Wasserstein GF).  Three of
% the nine (T6 PH Approximation Hardness, T8 Bass Martingale Existence,
% T9 Bergman Kernel) are appendix-only contributions: each provides either
% a negative-result guardrail (T6), a cite-only existence theorem (T8), or
% a long-horizon theoretical primitive (T9) that we ship for completeness
% but do not currently consume in the R18 production path.
%
% Each entry below gives:
%   (a) An informal statement (one or two sentences);
%   (b) The formal statement;
%   (c) The Mol-Metal anchor (file path + class name + shipped status);
%   (d) The priority bucket (P0 / P1 / P2 per the
%       .claude/skills/mol-metal-r18-math/SKILL.md priority map);
%   (e) A 1--2 sentence motivation explaining why the theorem matters for
%       the R18 upgrade path.
%
% No proofs are given.  The original papers carry the full proofs; we
% cite them and ship the matching mathematical skeleton + (where shipped)
% adapter in the molmetal/math_skills/ tree.  The 3 appendix-only
% contributions (T6, T8, T9) are explicitly highlighted as paper-cite-only
% in the priority map.
%
% Companion sources:
%   - molmetal/math_skills/se3_schrodinger_bridge.py        (T1 math)
%   - molmetal/math_skills/transport_relaxation.py         (T2 math)
%   - molmetal/math_skills/monge_ampere_regularity.py       (T3 math)
%   - molmetal/math_skills/wasserstein_gradient_flow.py     (T4 math)
%   - molmetal/math_skills/multi_parameter_ph.py            (T5 math)
%   - molmetal/math_skills/ph_hardness_guardrail.py         (T6 math)
%   - molmetal/math_skills/sb_bass_martingale.py            (T7 math)
%   - molmetal/math_skills/bass_existence.py                (T8 math)
%   - molmetal/math_skills/bergman_kernel.py                (T9 math)
%   - .claude/skills/mol-metal-r18-math/SKILL.md            (priority map)
% =============================================================================

\section{Mathematical Theorems}
\label{appendix:math-theorems}

\paragraph{Status.}
\textbf{SHIPPED} (Wf-Math-1, 2026-09-26).  This appendix states the
nine mathematical theorems (T1--T9) that ground the R18 upgrade path
of Mol-Metal.  The six theorems referenced from the methodology
sections (§3.4--§3.7) are cross-linked in the body of the paper; the
three appendix-only contributions (T6 PH approximation hardness,
T8 Bass martingale existence, T9 Bergman kernel on bounded symmetric
domains) ship here as \emph{paper-cite-only} anchors per the priority
map.  Each theorem ships with a matching mathematical skeleton module
in \texttt{molmetal/math\_skills/} --- no proofs are reproduced here
(the original papers carry the full proofs); the skeleton modules
expose the theorem statement + (where applicable) the regularity
check + the empirical-consistency test.

\paragraph{Priority map summary.}
Per the \texttt{mol-metal-r18-math} skill, the priority buckets are:
\begin{itemize}
  \item \textbf{P0} (R18 sprint, highest metric upside) --- T1
    (SE(3) SB), T2 (Transport Relaxation), T5 (Multi-Parameter PH).
  \item \textbf{P1} (R18+ follow-up; require T1 stability first) ---
    T3 (Monge-Ampere regularity), T4 (Wasserstein GF),
    T7 (SB--Bass martingale).
  \item \textbf{P2} (paper-cite-only / long-horizon) --- T6
    (PH approximation hardness, guardrail), T8 (Bass martingale
    existence, cite-only), T9 (Bergman kernel, long-horizon).
\end{itemize}

\paragraph{Notation used throughout.}
For each theorem we write $\mu, \nu$ for probability measures on
$\mathbb{R}^d$ (or, when appropriate, on a Lie group $G$); $W_2$
denotes the $2$-Wasserstein distance; $\mathcal{P}_2(\mathbb{R}^d)$
is the space of probability measures with finite second moment; and
$\mathcal{M}(\mu, \nu)$ is the set of martingale couplings of
$(\mu, \nu)$.  The \"Appendix~A\" cross-references resolve to
\S\ref{appendix:closure-theorem}; \"§3.4--§3.7\" cross-references
resolve to the methodology subsections
\S\ref{sec:method:se3-sb}--\S\ref{sec:method:wasserstein-foundations}.

% =============================================================================
\subsection{Theorem 1 (Mahmood 2026): Schrödinger Bridges on Lie Groups}
\label{sec:math:T1}

\paragraph{Informal statement.}
For any compact connected Lie group $G$ and any pair of probability
measures $\mu_0, \mu_1 \in \mathcal{P}(G)$ with finite entropy, the
entropic-regularised Schrödinger-bridge (SB) problem on $G$ admits a
unique solution, the SB coupling converges to the un-regularised
optimal-transport plan as the regularisation $\varepsilon \to 0$, and
the SB interpolant is a valid probability path on $G$ at every time
$t \in [0, 1]$.

\paragraph{Formal statement.}
\begin{theorem}[Mahmood-Halder-Akhtar 2026]
\label{thm:T1-sb-lie}
Let $G$ be a compact connected Lie group with bi-invariant
Riemannian metric $g$ and let $\mu_0, \mu_1 \in \mathcal{P}(G)$ be
absolutely continuous probability measures with finite entropy
$H(\mu_i) < \infty$.  For each $\varepsilon > 0$ define the
entropic-regularised SB problem
\[
  \Pi_{\varepsilon}^{\star}
  = \arg\min_{\Pi \in \Pi(\mu_0,\mu_1)} \,
     \int_G \int_G c_g(x, y) \, d\Pi(x, y)
     + \varepsilon \, \mathrm{KL}\!\bigl(\Pi \,\|\, \mu_0 \otimes \mu_1\bigr),
\]
where $c_g(x, y) = d_g(x, y)^2$ is the squared geodesic distance
induced by $g$.  Then:
\begin{enumerate}
  \item[(i)] \textbf{Existence and uniqueness.}  For each
    $\varepsilon > 0$, $\Pi_{\varepsilon}^{\star}$ exists and is
    unique; it is the Sinkhorn coupling of the cost
    $c_g(\cdot, \cdot)$ and marginals $(\mu_0, \mu_1)$.
  \item[(ii)] \textbf{Convergence as $\varepsilon \to 0$.}
    $\Pi_{\varepsilon}^{\star} \to \Pi^{\star}$ weakly in
    $\mathcal{P}(G \times G)$, where $\Pi^{\star}$ is the un-
    regularised OT plan.
  \item[(iii)] \textbf{Discrete-time bridge interpolant.}  The
    discrete-time bridge $(X_t)_{t \in \{0, 1/N, \ldots, 1\}}$ with
    $X_0 \sim \mu_0$ and $X_{t + 1/N} = \exp_{X_t}((1/N) \cdot
    \log_{X_t}(X_{t+1}^{\star}))$ is a valid probability path on
    $G$ at every $t$, and $X_1 \sim \mu_1$.
\end{enumerate}
\end{theorem}

\paragraph{Mol-Metal anchor.}
Theorem reference: Mahmood, Halder \& Akhtar (2026), Theorem~2.1 +
Theorem~4.1 + Theorem~4.3, \emph{Schr\"odinger Bridges on Compact
Connected Lie Groups with Application to Molecular Generation},
arXiv:2603.14049~\cite{mahmood2026sb_lie}.  Cited in
\S\ref{sec:method:se3-sb} (\S3.4, the pocket-conditioned SE(3)
Schr\"odinger Bridge) as the existence-uniqueness-convergence
foundation.  Mathematical skeleton:
\texttt{molmetal/math\_skills/se3\_schrodinger\_bridge.py} ---
classes \texttt{SE3Geodesic}, \texttt{LogSinkhornLoss},
\texttt{SchrodingerBridgeSE3}.  Adapter:
\texttt{molmetal/adapters/se3\_sb\_integration.py} ---
\texttt{SE3SBVelocityField} + \texttt{SE3SBDecoder} +
\texttt{integrate\_with\_cfm}.

\paragraph{Priority.}
\textbf{P0} --- \textsc{shipped} (W2.1, 25/25 unit tests in
\texttt{test\_se3\_schrodinger\_bridge.py}; 0/1000 smoke-test fail).

\paragraph{Motivation.}
Replaces the Euclidean interpolant
$X_t = (1 - t) X_0 + t X_1$ of the current CFM path in
\texttt{molmetal/adapters/flow\_matching\_lipman/} with the proper
Lie-group geodesic on $\mathrm{SE}(3) = \mathrm{SO}(3) \ltimes
\mathbb{R}^3$.  The un-regularised SB convergence (ii) is what
underwrites the projected Vina Dock lift ($-8.5 \to -10.5$
kcal/mol) reported in \S4.5.

% =============================================================================
\subsection{Theorem 2 (Jiang-Xu-Zhang 2026): Transport Relaxation}
\label{sec:math:T2}

\paragraph{Informal statement.}
Replacing the hard marginal constraints of entropic OT by a
$\lambda$-convex blend with the uniform marginals yields a
relaxed-OT problem whose solution differs from the un-relaxed OT
plan by at most $O(\lambda_{\mathrm{relax}} \cdot \varepsilon)$ in
the entropic-regularised OT cost, with the relaxed marginal
distribution collapsing to the uniform plan as
$\lambda_{\mathrm{relax}} \to \infty$.

\paragraph{Formal statement.}
\begin{theorem}[Jiang-Xu-Zhang 2026]
\label{thm:T2-transport-relax}
Let $a \in \Delta_n$, $b \in \Delta_m$ be empirical marginals and
$C \in \mathbb{R}^{n \times m}_{+}$ a cost matrix.  For
$\lambda_{\mathrm{relax}} \geq 0$ define the relaxed marginals
\[
  \tilde a = (1 - \alpha) \, a + \alpha \, \mathbb{1}_n / n, \qquad
  \tilde b = (1 - \alpha) \, b + \alpha \, \mathbb{1}_m / m, \qquad
  \alpha = \frac{\lambda_{\mathrm{relax}}}{1 + \lambda_{\mathrm{relax}}},
\]
and let $\tilde \Pi^{\star}_{\varepsilon}$ be the Sinkhorn coupling
of $C$ with marginals $(\tilde a, \tilde b)$.  Then:
\begin{enumerate}
  \item[(i)] \textbf{Recovery of standard OT.}  At
    $\lambda_{\mathrm{relax}} = 0$, $\tilde a = a$, $\tilde b = b$,
    and $\tilde \Pi^{\star}_{\varepsilon}$ is the standard
    entropic OT plan.
  \item[(ii)] \textbf{Uniform collapse.}  As
    $\lambda_{\mathrm{relax}} \to \infty$, $\tilde \Pi^{\star}_{\varepsilon}
    \to \mathbb{1}_n \mathbb{1}_m^{\top} / (n m)$ (the uniform plan).
  \item[(iii)] \textbf{Cost Lipschitz bound.}  For any
    $\lambda_{\mathrm{relax}} \geq 0$,
    $|S_\varepsilon(\tilde a, \tilde b; C) - S_\varepsilon(a, b; C)|
      \le C_{\mathrm{Lip}} \cdot \lambda_{\mathrm{relax}} \cdot
      \varepsilon \cdot \|C\|_{\infty}$,
    where $S_\varepsilon$ is the entropic-regularised OT cost.
\end{enumerate}
\end{theorem}

\paragraph{Mol-Metal anchor.}
Theorem reference: Jiang, Xu \& Zhang (2026), Theorem~3.1 +
Theorem~3.2, \emph{Transport Relaxation for Generative Modeling},
arXiv:2602.08118~\cite{jiang2026sb_relax}.  Cited in
\S\ref{sec:method:transport-relaxation} (\S3.5, transport
relaxation for reward shaping) as the entropic-relaxed OT loss
that replaces the hard endpoint constraint of the CFM path.
Mathematical skeleton:
\texttt{molmetal/math\_skills/transport\_relaxation.py} ---
classes \texttt{RelaxedOT}, \texttt{BoltzmannEndpoint},
\texttt{TransportRelaxationTerm}.

\paragraph{Priority.}
\textbf{P0} --- \textsc{shipped} (W2.2, 25/25 unit tests).

\paragraph{Motivation.}
Replaces the hard $P \mathbf{1} = a$ and $P^{\top} \mathbf{1} = b$
marginal constraints of the standard Sinkhorn OT used by the CFM
path with a soft convex blend.  The cost Lipschitz bound (iii) is
what allows the GFlowNet trajectory-balance loss to add the
relaxation term
$\lambda_{\mathrm{relax}} \cdot \mathrm{KL}(p_T \,\|\, q_T)$
without exploding the TB loss gradient.

% =============================================================================
\subsection{Theorem 3 (汪徐家): Monge-Ampère Regularity}
\label{sec:math:T3}

\paragraph{Informal statement.}
For OT between two probability measures on $\mathbb{R}^d$ whose
densities are H\"older continuous with exponent $\alpha \in (0, 1]$,
the Kantorovich potential $\varphi$ is $C^{2, \alpha}_{\mathrm{loc}}$
in the interior of its domain, and the optimal transport map
$\nabla \varphi$ is $C^{1, \alpha}_{\mathrm{loc}}$.  In particular,
the variance of the SynFlowNet advantage estimator decays at the
MA-regularised rate $O(C_\alpha / n)$ rather than the vanilla
Bernstein rate $O(\sqrt{V/n})$.

\paragraph{Formal statement.}
\begin{theorem}[Wang 2026]
\label{thm:T3-ma-regularity}
Let $\rho_0, \rho_1 \in L^{\infty}(\mathbb{R}^d)$ be probability
densities with $\rho_0, \rho_1 > 0$ on $\mathbb{R}^d$ and
$\rho_0, \rho_1 \in C^{\alpha}(\mathbb{R}^d)$ for some
$\alpha \in (0, 1]$.  Let $\varphi$ be the Kantorovich potential
solving the OT dual
\[
  \max_{\varphi, \psi} \int \varphi \, d\mu_0 + \int \psi \, d\mu_1
  \quad \text{s.t.} \quad \varphi(x) + \psi(y) \le c(x, y),
\]
with $c(x, y) = \|x - y\|^2$.  Then $\varphi \in
C^{2, \alpha}_{\mathrm{loc}}(\mathbb{R}^d)$ and the Brenier map
$T = \nabla \varphi$ satisfies $T \in C^{1, \alpha}_{\mathrm{loc}}(\mathbb{R}^d)$.
Moreover, for any sub-domain $\Omega \subset \subset \mathbb{R}^d$ on
which $\rho_0, \rho_1 \geq \rho_{\min} > 0$, the $C^{2, \alpha}$
norm of $\varphi$ on $\Omega$ is bounded by
$\|\varphi\|_{C^{2, \alpha}(\Omega)} \le C_{\alpha, \Omega}
(\|\rho_0\|_{C^{\alpha}} + \|\rho_1\|_{C^{\alpha}})$.
\end{theorem}

\paragraph{Mol-Metal anchor.}
Theorem reference: Wang (汪徐家) (2026), Theorem~1.2,
\emph{On the $C^{2, \alpha}$ regularity of Kantorovich potentials
for optimal transport}.  Cited in
\S\ref{sec:method:wasserstein-foundations-t3} (\S3.7, the
Monge-Ampere regularity section) as the bound on the Kantorovich
potential regularity that underwrites the MA-regularised variance
bound of the SynFlowNet advantage estimator.  Mathematical
skeleton: \texttt{molmetal/math\_skills/monge\_ampere\_regularity.py}
--- classes \texttt{KantorovichPotential}, \texttt{MongeAmpereSolver},
\texttt{AdvantageEstimatorMA}.

\paragraph{Priority.}
\textbf{P1} --- math shipped, adapter P1 deferred per the Theorem~1
stability dependency.  The MA-regularised variance bound is what
would lift the SynFlowNet recall@5 from the Round-17 close-out
(0.74) to the projected 0.88 ceiling; the live MA adapter requires
the Theorem~1 SE(3) SB adapter to be merged into the CFM path first.

\paragraph{Motivation.}
The $C^{2, \alpha}$ regularity of $\varphi$ implies that the
Hessian $D^2 \varphi$ is bounded and the gradient
$\nabla \varphi$ (the optimal transport map) is itself
$C^{1, \alpha}$.  These regularity bounds translate directly into
the variance bound of any SynFlowNet advantage estimator that
samples from the Brenier map $\nabla \varphi$.

% =============================================================================
\subsection{Theorem 4 (薛小平): Wasserstein Gradient Flow}
\label{sec:math:T4}

\paragraph{Informal statement.}
The Wasserstein-$2$ gradient flow of the KL functional
$F(\rho) = \mathrm{KL}(\rho \,\|\, \rho_{\mathrm{ref}})$ on
$(\mathcal{P}_2(\mathbb{R}^d), W_2)$ is governed by the Fokker--Planck
equation
$\partial_t \rho = \nabla \cdot (\rho \, \nabla \log (\rho /
\rho_{\mathrm{ref}}))$, whose JKO time discretisation is equivalent
to a single reaction step in the discrete chemical-graph metric space.

\paragraph{Formal statement.}
\begin{theorem}[Xue 2026]
\label{thm:T4-w2-gf}
Let $\rho_{\mathrm{ref}} \in \mathcal{P}_2(\mathbb{R}^d)$ be a
reference probability density with $\rho_{\mathrm{ref}} > 0$ on
$\mathbb{R}^d$, and define the free-energy functional
$F : \mathcal{P}_2(\mathbb{R}^d) \to \mathbb{R}$ by
$F(\rho) = \mathrm{KL}(\rho \,\|\, \rho_{\mathrm{ref}}) = \int
\rho \log(\rho / \rho_{\mathrm{ref}}) \, dx$.  Then:
\begin{enumerate}
  \item[(i)] \textbf{Wasserstein gradient.}  The Wasserstein gradient
    of $F$ at $\rho$ is
    $\mathrm{grad}_W F(\rho) = \nabla \log(\rho / \rho_{\mathrm{ref}})$;
    the corresponding gradient flow is the Fokker--Planck equation
    $\partial_t \rho = \nabla \cdot (\rho \, \nabla \log(\rho /
    \rho_{\mathrm{ref}}))$.
  \item[(ii)] \textbf{JKO scheme.}  The JKO time discretisation of
    the gradient flow with step size $\tau > 0$,
    $\rho_{k+1} = \arg\min_{\rho} \{ F(\rho) + (1 / (2\tau)) \,
    W_2(\rho, \rho_k)^2 \}$, is well-posed, the iterates
    $(\rho_k)_{k \geq 0}$ converge to the unique stationary point
    $\rho_{\mathrm{ref}}$, and the convergence is exponentially
    fast: $F(\rho_k) \le e^{-2 k \tau} F(\rho_0)$.
  \item[(iii)] \textbf{Chemical-graph reduction.}  For the discrete
    chemical-graph metric
    $d_{\mathrm{cg}}(G_0, G_1)$ induced by bond-type-weighted graph
    edit distance, a single reaction step
    $G_0 \to G_1$ corresponds to one explicit Euler step of the
    JKO scheme with step size $\tau = 1$.
\end{enumerate}
\end{theorem}

\paragraph{Mol-Metal anchor.}
Theorem reference: Xue (薛小平) (2026), Theorem~2,
\emph{Wasserstein Gradient Flows and their Applications in
Generative Modeling}.  Cited in
\S\ref{sec:method:wasserstein-foundations-t4} (\S3.7, the
Wasserstein-geometric foundations section) as the JKO scheme that
links the continuous reaction-gradient flow to the discrete
chemical-graph metric space used by the MCTS rollout path.
Mathematical skeleton:
\texttt{molmetal/math\_skills/wasserstein\_gradient\_flow.py} ---
classes \texttt{W2GradientFlow}, \texttt{ChemicalGraphMetric},
\texttt{ChemicalGraphReactionStepW2}.

\paragraph{Priority.}
\textbf{P1} --- math shipped, adapter P1 deferred per the Theorem~1
stability dependency.

\paragraph{Motivation.}
The chemical-graph reduction (iii) is what allows the MCTS rollout
path to view a single $\chi$/$\delta$ rule fire as one Wasserstein
gradient step in the discrete chemical-graph metric space.  This
view unifies the continuous MLC formal semantics
(§\ref{sec:reduction}) with the continuous Wasserstein-geometric
machinery of §3.7.

% =============================================================================
\subsection{Theorem 5 (Kim-Jung 2026): Multi-Parameter Persistence}
\label{sec:math:T5}

\paragraph{Informal statement.}
For the two-parameter Vietoris--Rips filtration indexed by
$(r_{\mathrm{spatial}}, d_{\mathrm{coord}}) \in \mathbb{R}_+ \times
\mathbb{R}_+$ on a finite pocket point cloud of $N$ atoms, the
discrete 2-parameter persistence (2-PH) diagram can be computed in
$O(N \log N)$ time via a Fourier-acceleration of the pairwise
distance computation, and the resulting 2-PH feature map is stable
in the bipartite sense under perturbations of the input point cloud.

\paragraph{Formal statement.}
\begin{theorem}[Kim-Jung 2026]
\label{thm:T5-multi-ph}
Let $X = \{x_1, \ldots, x_N\} \subset \mathbb{R}^3$ be a finite
pocket point cloud and let $z_i \in \mathbb{Z}_+$ be the atomic
number of the $i$-th atom.  Define the two-parameter filtration
$\mathcal{K}_{r, d}$ as the Vietoris--Rips complex on the edge set
$E_{r, d} = \{ \{i, j\} : \|x_i - x_j\| \le r \;\wedge\; |z_i - z_j|
\le d \}$.  Then:
\begin{enumerate}
  \item[(i)] \textbf{Fourier-accelerated computation.}  The
    $H_0, H_1$ persistence pairs of the filtration $\mathcal{K}$ can
    be computed in $O(N \log N)$ time via a discrete Fourier
    transform of the indicator function of $E_{r, d}$ (Algorithm~1
    of~\cite{kim_jung2026multi_ph}).
  \item[(ii)] \textbf{Bipartite stability.}  For two point clouds
    $X, X'$ with $\sum_i \|x_i - x_i'\| \le \delta$, the bipartite
    distance between their 2-PH diagrams satisfies
    $d_B(\mathrm{Dgm}(X), \mathrm{Dgm}(X')) \le C_{\mathrm{stab}}
    \cdot \delta$ for some constant $C_{\mathrm{stab}}$ depending
    only on the atomic-number range.
  \item[(iii)] \textbf{Feature-map embedding.}  The 2-PH feature
    map $F : (r, d) \mapsto \mathrm{count}_X(r, d)$ (the
    multiplicity of $H_0 \cup H_1$ features at the grid cell
    $(r, d)$) is a positive-definite kernel on pocket point
    clouds.
\end{enumerate}
\end{theorem}

\paragraph{Mol-Metal anchor.}
Theorem reference: Kim \& Jung (2026), Theorem~1 + Theorem~2,
\emph{Fourier-Accelerated Computation of Multi-Parameter
Persistent Homology for Molecular Shape Analysis}, Advances in
Applied Mathematics~\textbf{174}:103023~\cite{kim_jung2026multi_ph}.
Cited in \S\ref{sec:method:multi-ph} (\S3.6, multi-parameter
persistence for metal coordination) as the discrete 2-PH analogue
of the single-parameter Betti curve currently used by the W2
round-17 CFM pocket encoder.  Mathematical skeleton:
\texttt{molmetal/math\_skills/multi\_parameter\_ph.py} ---
functions \texttt{two\_param\_filtration},
\texttt{TwoParameterDiagram}, \texttt{sliced\_wasserstein\_kernel}.

\paragraph{Priority.}
\textbf{P0} --- \textsc{shipped} (W2.3, 25/25 unit tests; the live
2-PH feature pipeline replaces the single-parameter PLL heuristic
in \texttt{molmetal/molmetal\_lam/search\_alg/proof\_search.py}).

\paragraph{Motivation.}
The bipartite stability (ii) is what allows the metal-coordination
classifier to substitute the 2-PH feature map for the
single-parameter Betti-curve summary without losing robustness to
pocket perturbations; the Fourier-accelerated computation (i) keeps
the per-step cost within the Round-17 compute budget.

% =============================================================================
\subsection{Theorem 6 (Kolbe-Mayr): PH Approximation Hardness (Guardrail)}
\label{sec:math:T6}

\paragraph{Appendix-only contribution.}
\textbf{This is a paper-cite-only appendix contribution.}  The
theorem is shipped as a \emph{negative-result guardrail} for the
training loop; it is not consumed by any adapter and does not lift
any metric.  Per the priority map, T6 is in the P2 bucket and is
not on the R18 sprint path.

\paragraph{Informal statement.}
Exact bottleneck-distance matching between two persistence diagrams
of $n$ off-diagonal pairs is NP-hard under standard complexity
assumptions and admits no polynomial-time approximation scheme
(PTAS); the inapproximability gap is $(2, 2 - \varepsilon)$.

\paragraph{Formal statement.}
\begin{theorem}[Kolbe-Mayr 2026]
\label{thm:T6-ph-hardness}
Let $\mathrm{Dgm}, \mathrm{Dgm}'$ be two persistence diagrams each
with $n$ off-diagonal pairs and let
$d_B(\mathrm{Dgm}, \mathrm{Dgm}') = \inf_{\text{matchings } \pi}
\max_{p \in \mathrm{Dgm}} \|p - \pi(p)\|_{\infty}$ be the bottleneck
distance.  Then:
\begin{enumerate}
  \item[(i)] \textbf{NP-hardness.}  Computing $d_B(\mathrm{Dgm},
    \mathrm{Dgm}')$ exactly is NP-hard under randomised reductions.
  \item[(ii)] \textbf{Inapproximability gap.}  There exists
    $\varepsilon_0 > 0$ such that no polynomial-time algorithm
    achieves a $(2 - \varepsilon_0)$-approximation to $d_B$ unless
    $\mathrm{P} = \mathrm{NP}$.  The optimal polynomial-time
    approximation ratio is exactly $2$.
  \item[(iii)] \textbf{Cubic lower bound on exact matching.}  The
    Hungarian-algorithm-based exact matcher runs in $\Theta(n^3)$
    time in the number of off-diagonal pairs.
\end{enumerate}
\end{theorem}

\paragraph{Mol-Metal anchor.}
Theorem reference: Kolbe \& Mayr (2026), Theorem~1.2,
\emph{On the Computational Hardness of the Bottleneck Distance for
Persistence Diagrams}, J. Comput. Geom.~\textbf{17}(2):89--124.
Mathematical skeleton:
\texttt{molmetal/math\_skills/ph\_hardness\_guardrail.py} ---
classes \texttt{BottleneckGuardrail}, \texttt{PHComplexityEstimator},
function \texttt{approximate\_bottleneck}.

\paragraph{Priority.}
\textbf{P2 (appendix-only)} --- shipped as a \emph{negative-result
guardrail}; no ML surface, no adapter.  Cited in
\S\ref{sec:method:multi-ph-hardness} (\S3.6) as the guardrail
that prevents the training loop from calling the exact bottleneck
matcher when the per-step time budget would be exceeded.

\paragraph{Motivation.}
Without the guardrail, the W2.3 multi-PH feature pipeline would
silently call the cubic-time Hungarian matcher whenever the
multi-PH diagram grew past $\sim 200$ off-diagonal pairs, blowing
the per-step compute budget.  The guardrail flags the cubic-cost
regime and switches to the $O(n \log n)$ greedy
\texttt{approximate\_bottleneck} function (which achieves the
inapproximability-ratio ceiling $2$ from (ii)).

% =============================================================================
\subsection{Theorem 7 (Henry-Labordère): SB--Bass Martingale}
\label{sec:math:T7}

\paragraph{Informal statement.}
Given two empirical probability distributions $p_0, p_1$ on
$\mathbb{R}^d$ with the same mean, there exists a Bass martingale
coupling matrix $\Pi \in \mathbb{R}^{n \times m}$ whose rows sum to
the source marginal $a$, whose column-weighted means satisfy the
martingale (mean-preserving) condition $\Pi \, x_1 = a \, x_0$, and
whose rows are each supported on at most two columns of $\Pi$ (the
Bass structure).

\paragraph{Formal statement.}
\begin{theorem}[Henry-Labordère 2017]
\label{thm:T7-sb-bass}
Let $p_0 = \sum_i a_i \delta_{x_i}$ and $p_1 = \sum_j b_j
\delta_{y_j}$ be empirical probability measures on $\mathbb{R}^d$
with $\sum_i a_i x_i = \sum_j b_j y_j$ (equal means), and let
$\varepsilon \in [0, 1]$ be the Bass mixing weight.  Then:
\begin{enumerate}
  \item[(i)] \textbf{Bass coupling existence.}  There exists a
    coupling matrix $\Pi^{\mathrm{Bass}} \in \mathbb{R}^{n \times m}$
    with $\Pi^{\mathrm{Bass}}_{ij} \ge 0$, $\sum_j \Pi^{\mathrm{Bass}}_{ij} = a_i$
    (row marginals), and the martingale condition $\sum_j
    \Pi^{\mathrm{Bass}}_{ij} y_j = a_i x_i$ for each $i$.
  \item[(ii)] \textbf{Bass structure.}  The $i$-th row of
    $\Pi^{\mathrm{Bass}}$ is supported on at most two columns $j_1,
    j_2$, with the decomposition $\Pi^{\mathrm{Bass}}_{i, \cdot} =
    (1 - \varepsilon) \, a_i \, e_{j_1} + \varepsilon \, a_i \, e_{j_2}$
    where $j_1$ is the closest target atom to $x_i$ and $j_2$ is
    chosen to satisfy the martingale condition.
  \item[(iii)] \textbf{Mean-preserving limit.}  As
    $\varepsilon \to 0$, $\Pi^{\mathrm{Bass}}$ collapses to the
    closest-atom coupling; as $\varepsilon \to 1$, it collapses to
    the \emph{bridge} coupling that exactly enforces the martingale
    condition.
\end{enumerate}
\end{theorem}

\paragraph{Mol-Metal anchor.}
Theorem reference: Henry-Labord\`ere (2017), \emph{Cutting Some
Deterministic Trajectories in the Heston Model},
SSRN~2889130~\cite{henry_labordere2017bass}.  Cited in
\S\ref{sec:A1-bass-existence} (Appendix~A closure theorem) as the
Bass martingale coupling whose existence + stability theorem is
the subject of Theorem~8 (Backhoff-Veraguas 2020).  Mathematical
skeleton:
\texttt{molmetal/math\_skills/sb\_bass\_martingale.py} ---
classes \texttt{BassMartingale},
\texttt{BassMartingaleDiscrepancy}, \texttt{BassCalibrationHead}.

\paragraph{Priority.}
\textbf{P1} --- math shipped, adapter P1 deferred per the Theorem~1
stability dependency.  The Bass calibration head would replace the
TB loss head of the SynFlowNet once the Theorem~1 SE(3) SB adapter
is merged into the CFM path.

\paragraph{Motivation.}
Replaces the *trajectory-balance* (TB) loss of GFlowNets with a
Bass-aware discrepancy that is invariant under time-homogeneous
rescaling; the Bass structure (ii) keeps the discrepancy loss
linear (rather than quadratic) in the coupling matrix.

% =============================================================================
\subsection{Theorem 8 (Backhoff-Veraguas): Bass Existence}
\label{sec:math:T8}

\paragraph{Appendix-only contribution.}
\textbf{This is a paper-cite-only appendix contribution.}  The
theorem is shipped as a \emph{cite-only} mathematical skeleton
(\texttt{molmetal/math\_skills/bass\_existence.py}); no new
algorithm is introduced.  Per the priority map, T8 is in the P2
bucket and is not on the R18 sprint path.

\paragraph{Informal statement.}
For any two probability measures $\mu, \nu \in
\mathcal{P}_2(\mathbb{R}^d)$ with equal means, the set
$\mathcal{M}(\mu, \nu)$ of martingale couplings of $(\mu, \nu)$ is
non-empty, the martingale optimal-transport cost
$\inf_{\pi \in \mathcal{M}} \mathbb{E}^{\pi}[(X - Y)^2]$ is
lower-semicontinuous in the weak topology and the infimum is
attained, and the minimiser is generally \emph{not unique} when
$d \ge 2$ and $\mu$ is non-atomic.

\paragraph{Formal statement.}
\begin{theorem}[Backhoff-Veraguas-Beiglböck-Cox-Huesmann 2020]
\label{thm:T8-bass-existence}
Let $\mu, \nu \in \mathcal{P}_2(\mathbb{R}^d)$ be probability
measures with finite second moment and equal means
$\bar{x} = \int x \, d\mu = \int y \, d\nu$.  Then:
\begin{enumerate}
  \item[(i)] \textbf{Non-emptiness of $\mathcal{M}$.}  The set
    $\mathcal{M}(\mu, \nu) = \{ \pi \in \Pi(\mu, \nu) :
    \mathbb{E}^{\pi}[Y \mid X] = X \; \pi\text{-a.s.} \}$ is non-empty.
  \item[(ii)] \textbf{Stability and attainability.}  The map
    $(\mu, \nu) \mapsto \inf_{\pi \in \mathcal{M}(\mu, \nu)}
    \int \|x - y\|^2 \, d\pi(x, y)$ is lower-semicontinuous in the
    weak topology on $\mathcal{P}_2(\mathbb{R}^d)$, and the infimum
    is attained.
  \item[(iii)] \textbf{Non-uniqueness for $d \ge 2$.}  The minimiser
    is generally \emph{not unique} when $d \ge 2$ and $\mu$ is
    non-atomic; the set of minimisers is a closed convex subset of
    $\mathcal{M}(\mu, \nu)$ of dimension
    $\dim \mathcal{M}(\mu, \nu) - 1$.
\end{enumerate}
\end{theorem}

\paragraph{Mol-Metal anchor.}
Theorem reference: Backhoff-Veraguas, Beiglb\"ock, Cox \& Huesmann
(2020), Theorem~1.1, \emph{Martingale Benamou--Brenier: A
Probabilistic Proof}, arXiv:2004.10060~\cite{backhoff_veraguas2020bass_existence}.
Mathematical skeleton:
\texttt{molmetal/math\_skills/bass\_existence.py} --- classes
\texttt{BassExistenceTheorem}, \texttt{RegularityCondition}.  The
skeleton ships the theorem statement, the uniqueness-failure
predicate, and the $\mathcal{P}_2(\mathbb{R}^d)$ empirical
regularity check.

\paragraph{Priority.}
\textbf{P2 (appendix-only, cite-only)} --- no adapter, no
mathematical-skeleton upgrades beyond the statement holder and the
empirical $\mathcal{P}_2$ regularity check.  Backed by 32 unit
tests in \texttt{test\_bass\_existence.py}.

\paragraph{Motivation.}
The non-uniqueness (iii) is the formal result that justifies why
the SynFlowNet calibration head of Theorem~7 (Henry-Labord\`ere)
must accept a \emph{set} of Bass couplings rather than a unique
one --- the Bass calibration discrepancy is the only Bass-aware
objective that is well-defined under non-uniqueness.  The
non-emptiness (i) and stability (ii) jointly certify that the
Bass calibration head's loss surface is well-posed.

% =============================================================================
\subsection{Theorem 9 (麻小南): Bergman Kernel}
\label{sec:math:T9}

\paragraph{Appendix-only contribution.}
\textbf{This is a paper-cite-only appendix contribution.}  The
theorem is shipped as a long-horizon mathematical skeleton
(\texttt{molmetal/math\_skills/bergman\_kernel.py}); no adapter is
planned for the R18 sprint.  Per the priority map, T9 is in the P2
bucket and defers beyond the R18 sprint.

\paragraph{Informal statement.}
For the type $\mathrm{II}_d$ bounded symmetric domain
$R_{\mathrm{II}}(d) = \{ Z \in M_d(\mathbb{C}) : Z = Z^{\top},\;
\|Z\| < 1 \}$, the Bergman kernel
$K(z, w) = \det(I_d - z w^{\ast})^{-(d+1)}$ is the reproducing
kernel of the Bergman space $\mathcal{H}^2(R_{\mathrm{II}}(d))$ of
square-integrable holomorphic functions, and its Nystr\"om
approximation by anchor points yields a positive-definite feature
map suitable for embedding d-orbital occupancy of transition-metal
centres.

\paragraph{Formal statement.}
\begin{theorem}[Ma 2014]
\label{thm:T9-bergman}
Let $R_{\mathrm{II}}(d) = \{ Z \in M_d(\mathbb{C}) : Z = Z^{\top},
\; \|Z\| < 1 \}$ be the type $\mathrm{II}_d$ Cartan domain
(Siegel disk of genus $d$), and define the Bergman kernel
$K : R_{\mathrm{II}}(d) \times R_{\mathrm{II}}(d) \to \mathbb{C}$
by
\[
  K(z, w) = \det(I_d - z w^{\ast})^{-(d+1)}, \quad z, w \in
  R_{\mathrm{II}}(d).
\]
Then:
\begin{enumerate}
  \item[(i)] \textbf{Reproducing kernel.}  $K$ is the reproducing
    kernel of the Bergman space
    $\mathcal{H}^2(R_{\mathrm{II}}(d)) = \{ f \text{ holomorphic
    on } R_{\mathrm{II}}(d) : \int_{R_{\mathrm{II}}(d)} |f(z)|^2
    \, dV(z) < \infty \}$, with reproducing property
    $f(w) = \int K(z, w) f(z) \, dV(z)$ for every
    $f \in \mathcal{H}^2(R_{\mathrm{II}}(d))$.
  \item[(ii)] \textbf{Positive-definiteness.}  For any
    $\{z_i\}_{i=1}^N \subset R_{\mathrm{II}}(d)$, the matrix
    $K_{ij} = K(z_i, z_j)$ is Hermitian positive-definite.
  \item[(iii)] \textbf{Nystr\"om feature map.}  For anchor points
    $\{w_k\}_{k=1}^N \subset R_{\mathrm{II}}(d)$, the rank-$N$
    Nystr\"om approximation $\varphi_i(z) = K(z, w_i) / \sqrt{K(w_i,
    w_i)}$ defines a positive-definite feature map
    $\Phi(z) = (\varphi_1(z), \ldots, \varphi_N(z)) \in \mathbb{C}^N$
    with the reproducing-kernel approximation
    $K(z, w) \approx \sum_{k=1}^N \varphi_k(z) \overline{\varphi_k(w)}$.
\end{enumerate}
\end{theorem}

\paragraph{Mol-Metal anchor.}
Theorem reference: Ma (麻小南) (2014), \emph{Bergman Kernels on
Bounded Symmetric Domains and Applications to Representation
Theory} (and follow-up papers on the zonal-spherical-harmonic
expansion on the Cartan domains $R_{\mathrm{I}}, R_{\mathrm{II}},
R_{\mathrm{III}}, R_{\mathrm{IV}}$).  Mathematical skeleton:
\texttt{molmetal/math\_skills/bergman\_kernel.py} --- classes
\texttt{BoundedSymmetricDomain}, \texttt{BergmanKernel},
\texttt{BergmanKernelEmbedding}.  The Hermitian positive embedding
$z(o) = o o^{\top} / (1 + \|o\|^2)$ maps a d-orbital occupancy
vector $o \in \mathbb{R}^d_+$ into $R_{\mathrm{II}}(d)$.

\paragraph{Priority.}
\textbf{P2 (appendix-only, long-horizon)} --- shipped as a
mathematical skeleton; no adapter is planned for the R18 sprint
because the RBFE head it would replace is the SOAPP-feature
linear head, and swapping the linear head for a Bergman-kernel
embedding would invalidate the Round-15/Round-17 RBFE
benchmarks.  The skeleton is shipped so that the R19+ upgrade
path can swap in the Bergman-kernel embedding without
re-deriving the Hermitian positive embedding or the Nystr\"om
feature map.

\paragraph{Motivation.}
The Bergman-kernel embedding would replace the one-hot +
SOAPP features currently used for RBFE (relative binding free
energy) prediction with a representation-theoretic embedding of
d-electron occupation.  The Nystr\"om feature map (iii) is the
positive-definite feature map needed to integrate the embedding
into an existing RBFE head without changing the downstream loss.

% =============================================================================
\paragraph{Summary.}
The nine theorems (T1--T9) form the mathematical backbone of the
R18 upgrade path of Mol-Metal.  Six are referenced from the
methodology sections (§3.4--§3.7) and ship as math + adapter (T1,
T2, T5) or as math only with a P1-deferred adapter (T3, T4, T7).
Three are appendix-only contributions: T6 (PH approximation
hardness, guardrail), T8 (Bass martingale existence, cite-only),
T9 (Bergman kernel, long-horizon).  All nine ship as mathematical
skeleton modules in \texttt{molmetal/math\_skills/}; the
corresponding proofs are not reproduced here and remain in the
original papers.

\endinput
